\section{Conclusiones}

La línea base conceptual D1 a D6 define una plataforma de apoyo a la prevención cardiovascular, no un sistema de diagnóstico autónomo. La comparación entre Cloud, Local e Híbrida es válida porque las tres alternativas conservan el mismo alcance, los mismos supuestos de usuarios, registros, operación, reentrenamiento y continuidad.

El stack propuesto es portable, mantenible y predominantemente open source. GCP se usa como escenario Cloud concreto, sin imponer servicios avanzados de ML antes de que su necesidad esté demostrada. La alternativa Local ofrece control directo, pero exige una operación y postura de seguridad sostenida. La alternativa Híbrida equilibra control de identidad local con elasticidad analítica en cloud, a costa de mayor integración y operación dual.

El enfoque adaptativo es consistente con la incertidumbre de los datos y del proyecto. El informe maestro declara el proyecto \textbf{viable condicionado} y selecciona Cloud/GCP como alternativa integral recomendada para el piloto, con TCO oficial de \$10.815.487 CLP a 36 meses. La selección está condicionada al cumplimiento verificable de CA-E01 a CA-E10, autorización y gobierno de datos aplicables, controles técnicos y de seguridad, aceptación de riesgos y validación local; no declara cumplimiento legal definitivo ni umbrales clínicos no acordados. Híbrida se mantiene como contingencia ante una no aprobación legal, institucional o de gobierno de datos de Cloud. Local es técnicamente viable, pero no se recomienda para este piloto por su mayor TCO y carga operacional.
